\chapter{Usage of Chatbots and Other Tools}\label{ch:chatbots}

ChatGPT and Codex (OpenAI) and Claude (Anthropic), all AI-based language models,
were used as supporting tools throughout this project for understanding technical
concepts, programming, and scientific writing. They were mainly used to understand
the theoretical foundations on which this project builds (\cref{ch:theory}):
polyconvexity and input-convex neural networks, periodic homogenization and the
Hill--Mandel condition, and the condensation of internal variables for a consistent
homogenized tangent. In addition, they were used to understand the NGSolve and
PyTorch interfaces that this project connects, in particular the
\code{IntegrationRuleSpace}, periodic and \code{NumberSpace} function spaces and the
automatic-differentiation interface of PyTorch (\cref{ch:methods}).

During implementation, these tools were used to discuss code structure, investigate
programming errors, and evaluate possible solutions. Specific attention was given to
the consistent transfer of stress and tangent between PyTorch and NGSolve at the
quadrature points, the work-conjugate homogenized stress including the particle
reactions (\cref{eq:hm}), the objective form of the reference normalization
(\cref{eq:norm}), and the globalization of the macroscopic Newton solver
(\cref{sec:macro-method}). The tools also assisted in preparing the PBS job scripts
for the TU~Freiberg cluster and in developing the verification test cases
(\cref{ch:verification}). All suggestions were critically reviewed against the
governing equations and the library documentation, and were tested before being
used in the project.

Furthermore, these tools helped improve the clarity of the report, correct
grammatical errors, and structure the chapters. Additional support was provided for
\LaTeX{} formatting, TikZ diagrams, the program flowchart (\cref{ch:flow}) and the
plotting scripts. All suggested content was carefully reviewed against the
implemented methods, their mathematical formulations, the stored results
(\cref{app:repro}) and the verification results (\cref{ch:verification}).
